\subsection{শিক্ষকের তথ্য থেকে নির্বাচন ও হিসাব}

নিচের সারণিতে MSc হলো \textbf{Master of Science} এবং MA হলো \textbf{Master of Arts} ডিগ্রির সংক্ষিপ্ত রূপ।

\begin{flushleft}\begin{minipage}{\linewidth}\raggedright

\small

\begin{tabular}{|l|l|l|l|l|}
\hline

\chSixHeader
TeacherID & Name & Subject & Qualification & Salary\\
\hline
T01 & Salma & ICT & MSc & 55000\\\hline
T02 & Rahat & Bangla & MA & 48000\\\hline
T03 & Farid & Physics & MSc & 62000\\\hline
T04 & Nabila & ICT & MSc & 58000\\\hline
T05 & Karim & English & MA & 45000\\
\hline

\end{tabular}

\captionof{table}{Teacher Table}

\end{minipage}\end{flushleft}

\noindent\begin{minipage}{\linewidth}
\chSixExample{৭.৩৭}{ICT বিষয়ের শিক্ষকদের তথ্য দেখাও।}

\noindent\begin{minipage}[c]{.53\linewidth}
\begin{verbatim}
SELECT TeacherID, Name, Salary
FROM Teacher
WHERE Subject='ICT';
\end{verbatim}
\end{minipage}\hfill\begin{minipage}[c]{.04\linewidth}\raggedright{\color{chSixTeal}$\Rightarrow$}\end{minipage}\hfill
\begin{minipage}[c]{.40\linewidth}\raggedright\footnotesize\setlength{\tabcolsep}{2pt}\textbf{ফল}\par\smallskip
\begin{tabular}{|l|l|l|}\hline\chSixHeader
TeacherID & Name & Salary\\\hline
T01 & Salma & 55000\\\hline
T04 & Nabila & 58000\\
\hline\end{tabular}
\end{minipage}
\end{minipage}\par\smallskip

\noindent\begin{minipage}{\linewidth}
\chSixExample{৭.৩৮}{যেসব শিক্ষকের Salary 50000-এর বেশি, তাঁদের তথ্য দেখাও।}

\noindent\begin{minipage}[c]{.53\linewidth}
\begin{verbatim}
SELECT Name, Subject, Salary
FROM Teacher
WHERE Salary>50000;
\end{verbatim}
\end{minipage}\hfill\begin{minipage}[c]{.04\linewidth}\raggedright{\color{chSixTeal}$\Rightarrow$}\end{minipage}\hfill
\begin{minipage}[c]{.40\linewidth}\raggedright\footnotesize\setlength{\tabcolsep}{2pt}\textbf{ফল}\par\smallskip
\begin{tabular}{|l|l|l|}\hline\chSixHeader
Name & Subject & Salary\\\hline
Salma & ICT & 55000\\\hline
Farid & Physics & 62000\\\hline
Nabila & ICT & 58000\\
\hline\end{tabular}
\end{minipage}
\end{minipage}\par\smallskip

\noindent\begin{minipage}{\linewidth}
\chSixExample{৭.৩৯}{Salary-এর সর্বোচ্চ মান বের করো।}

\noindent\begin{minipage}[c]{.53\linewidth}
\begin{verbatim}
SELECT MAX(Salary) AS HighestSalary
FROM Teacher;
\end{verbatim}
\end{minipage}\hfill\begin{minipage}[c]{.04\linewidth}\raggedright{\color{chSixTeal}$\Rightarrow$}\end{minipage}\hfill
\begin{minipage}[c]{.40\linewidth}\raggedright\footnotesize\setlength{\tabcolsep}{2pt}\textbf{ফল}\par\smallskip
\begin{tabular}{|l|}\hline\chSixHeader
HighestSalary\\\hline
62000\\
\hline\end{tabular}
\end{minipage}
\end{minipage}\par\smallskip

\noindent\begin{minipage}{\linewidth}
\chSixExample{৭.৪০}{শিক্ষকদের গড় Salary বের করো।}

\noindent\begin{minipage}[c]{.53\linewidth}
\begin{verbatim}
SELECT AVG(Salary) AS AverageSalary
FROM Teacher;
\end{verbatim}
\end{minipage}\hfill\begin{minipage}[c]{.04\linewidth}\raggedright{\color{chSixTeal}$\Rightarrow$}\end{minipage}\hfill
\begin{minipage}[c]{.40\linewidth}\raggedright\footnotesize\setlength{\tabcolsep}{2pt}\textbf{ফল}\par\smallskip
\begin{tabular}{|l|}\hline\chSixHeader
AverageSalary\\\hline
53600\\
\hline\end{tabular}
\end{minipage}
\end{minipage}\par\smallskip

\subsection{খেলোয়াড়ের তথ্য বিশ্লেষণ}

\begin{flushleft}\begin{minipage}{\linewidth}\raggedright

\small

\begin{tabular}{|l|l|l|l|l|l|}
\hline

\chSixHeader
PlayerID & Name & Team & Game & Age & Score\\
\hline
P01 & Sakib & Tigers & Cricket & 24 & 135\\\hline
P02 & Rafi & Eagles & Football & 22 & 18\\\hline
P03 & Tamim & Tigers & Cricket & 25 & 98\\\hline
P04 & Nabil & Warriors & Cricket & 21 & 152\\\hline
P05 & Rima & Queens & Badminton & 20 & 75\\
\hline

\end{tabular}

\captionof{table}{Player Table}

\end{minipage}\end{flushleft}

\noindent\begin{minipage}{\linewidth}
\chSixExample{৭.৪১}{Cricket খেলোয়াড়দের তথ্য দেখাও।}

\noindent\begin{minipage}[c]{.53\linewidth}
\begin{verbatim}
SELECT PlayerID, Name, Score
FROM Player
WHERE Game='Cricket';
\end{verbatim}
\end{minipage}\hfill\begin{minipage}[c]{.04\linewidth}\raggedright{\color{chSixTeal}$\Rightarrow$}\end{minipage}\hfill
\begin{minipage}[c]{.40\linewidth}\raggedright\footnotesize\setlength{\tabcolsep}{2pt}\textbf{ফল}\par\smallskip
\begin{tabular}{|l|l|l|}\hline\chSixHeader
PlayerID & Name & Score\\\hline
P01 & Sakib & 135\\\hline
P03 & Tamim & 98\\\hline
P04 & Nabil & 152\\
\hline\end{tabular}
\end{minipage}
\end{minipage}\par\smallskip

\noindent\begin{minipage}{\linewidth}
\chSixExample{৭.৪২}{যেসব খেলোয়াড়ের Score 100-এর বেশি, তাঁদের তথ্য দেখাও।}

\noindent\begin{minipage}[c]{.53\linewidth}
\begin{verbatim}
SELECT Name, Score
FROM Player
WHERE Score>100;
\end{verbatim}
\end{minipage}\hfill\begin{minipage}[c]{.04\linewidth}\raggedright{\color{chSixTeal}$\Rightarrow$}\end{minipage}\hfill
\begin{minipage}[c]{.40\linewidth}\raggedright\footnotesize\setlength{\tabcolsep}{2pt}\textbf{ফল}\par\smallskip
\begin{tabular}{|l|l|}\hline\chSixHeader
Name & Score\\\hline
Sakib & 135\\\hline
Nabil & 152\\
\hline\end{tabular}
\end{minipage}
\end{minipage}\par\smallskip

\noindent\begin{minipage}{\linewidth}
\chSixExample{৭.৪৩}{Score-এর সর্বোচ্চ মান বের করো।}

\noindent\begin{minipage}[c]{.53\linewidth}
\begin{verbatim}
SELECT MAX(Score) AS HighestScore
FROM Player;
\end{verbatim}
\end{minipage}\hfill\begin{minipage}[c]{.04\linewidth}\raggedright{\color{chSixTeal}$\Rightarrow$}\end{minipage}\hfill
\begin{minipage}[c]{.40\linewidth}\raggedright\footnotesize\setlength{\tabcolsep}{2pt}\textbf{ফল}\par\smallskip
\begin{tabular}{|l|}\hline\chSixHeader
HighestScore\\\hline
152\\
\hline\end{tabular}
\end{minipage}
\end{minipage}\par\smallskip

\noindent\begin{minipage}{\linewidth}
\chSixExample{৭.৪৪}{Score-এর গড় বের করো।}

\noindent\begin{minipage}[c]{.53\linewidth}
\begin{verbatim}
SELECT AVG(Score) AS AverageScore
FROM Player;
\end{verbatim}
\end{minipage}\hfill\begin{minipage}[c]{.04\linewidth}\raggedright{\color{chSixTeal}$\Rightarrow$}\end{minipage}\hfill
\begin{minipage}[c]{.40\linewidth}\raggedright\footnotesize\setlength{\tabcolsep}{2pt}\textbf{ফল}\par\smallskip
\begin{tabular}{|l|}\hline\chSixHeader
AverageScore\\\hline
95.60\\
\hline\end{tabular}
\end{minipage}
\end{minipage}\par\smallskip

\subsection{শহর অনুযায়ী গ্রাহকের তথ্য বিশ্লেষণ}

\begin{flushleft}\begin{minipage}{\linewidth}\raggedright

\small

\begin{tabular}{|l|l|l|l|}
\hline

\chSixHeader
CustomerID & Name & Phone & City\\
\hline
C01 & Asif & 01710000001 & Dhaka\\\hline
C02 & Nila & 01710000002 & Rajshahi\\\hline
C03 & Anika & 01710000003 & Dhaka\\\hline
C04 & Karim & 01710000004 & Khulna\\\hline
C05 & Arif & 01710000005 & Dhaka\\\hline
C06 & Mitu & 01710000006 & Rajshahi\\
\hline

\end{tabular}

\captionof{table}{Customer Table}

\end{minipage}\end{flushleft}

\noindent\begin{minipage}{\linewidth}
\chSixExample{৭.৪৫}{Dhaka শহরের গ্রাহকদের তথ্য দেখাও।}

\noindent\begin{minipage}[c]{.53\linewidth}
\begin{verbatim}
SELECT CustomerID, Name
FROM Customer
WHERE City='Dhaka';
\end{verbatim}
\end{minipage}\hfill\begin{minipage}[c]{.04\linewidth}\raggedright{\color{chSixTeal}$\Rightarrow$}\end{minipage}\hfill
\begin{minipage}[c]{.40\linewidth}\raggedright\footnotesize\setlength{\tabcolsep}{2pt}\textbf{ফল}\par\smallskip
\begin{tabular}{|l|l|}\hline\chSixHeader
CustomerID & Name\\\hline
C01 & Asif\\\hline
C03 & Anika\\\hline
C05 & Arif\\
\hline\end{tabular}
\end{minipage}
\end{minipage}\par\smallskip

\noindent\begin{minipage}{\linewidth}
\chSixExample{৭.৪৬}{মোট গ্রাহকসংখ্যা বের করো।}

\noindent\begin{minipage}[c]{.53\linewidth}
\begin{verbatim}
SELECT COUNT(*) AS TotalCustomer
FROM Customer;
\end{verbatim}
\end{minipage}\hfill\begin{minipage}[c]{.04\linewidth}\raggedright{\color{chSixTeal}$\Rightarrow$}\end{minipage}\hfill
\begin{minipage}[c]{.40\linewidth}\raggedright\footnotesize\setlength{\tabcolsep}{2pt}\textbf{ফল}\par\smallskip
\begin{tabular}{|l|}\hline\chSixHeader
TotalCustomer\\\hline
6\\
\hline\end{tabular}
\end{minipage}
\end{minipage}\par\smallskip

\noindent\begin{minipage}{\linewidth}
\chSixExample{৭.৪৭}{যেসব গ্রাহকের Name A দিয়ে শুরু, তাঁদের তথ্য দেখাও।}

\noindent\begin{minipage}[c]{.53\linewidth}
\begin{verbatim}
SELECT CustomerID, Name
FROM Customer
WHERE Name LIKE 'A%';
\end{verbatim}
\end{minipage}\hfill\begin{minipage}[c]{.04\linewidth}\raggedright{\color{chSixTeal}$\Rightarrow$}\end{minipage}\hfill
\begin{minipage}[c]{.40\linewidth}\raggedright\footnotesize\setlength{\tabcolsep}{2pt}\textbf{ফল}\par\smallskip
\begin{tabular}{|l|l|}\hline\chSixHeader
CustomerID & Name\\\hline
C01 & Asif\\\hline
C03 & Anika\\\hline
C05 & Arif\\
\hline\end{tabular}
\end{minipage}
\end{minipage}\par\smallskip

\noindent\begin{minipage}{\linewidth}
\chSixExample{৭.৪৮}{যেসব City-তে একাধিক গ্রাহক আছেন, সেগুলো দেখাও।}

\noindent\begin{minipage}[c]{.53\linewidth}
\begin{verbatim}
SELECT City,
       COUNT(*) AS TotalCustomer
FROM Customer
GROUP BY City
HAVING COUNT(*)>1;
\end{verbatim}
\end{minipage}\hfill\begin{minipage}[c]{.04\linewidth}\raggedright{\color{chSixTeal}$\Rightarrow$}\end{minipage}\hfill
\begin{minipage}[c]{.40\linewidth}\raggedright\footnotesize\setlength{\tabcolsep}{2pt}\textbf{ফল}\par\smallskip
\begin{tabular}{|l|l|}\hline\chSixHeader
City & TotalCustomer\\\hline
Dhaka & 3\\\hline
Rajshahi & 2\\
\hline\end{tabular}
\end{minipage}
\end{minipage}\par\smallskip

\subsection{ব্যাংকের সম্পর্কিত টেবিলে SQL-এর প্রয়োগ}

ব্যাংকের ডেটাবেজ একটি গুরুত্বপূর্ণ বাস্তব উদাহরণ। এখানে একই সঙ্গে
শর্ত, সমষ্টিগত ফাংশন, UPDATE এবং সম্পর্কিত টেবিলের কুয়েরি
বোঝা যায়।

\begin{chSixSideBySide}[.48]
\subsubsection{BankCustomer Table}



\small

\begin{tabular}{|l|l|l|l|}
\hline

\chSixHeader
CustomerID & Name & Address & Phone\\
\hline
C01 & Rafi & Rajshahi & 01712000001\\\hline
C02 & Nila & Dhaka & 01712000002\\\hline
C03 & Sami & Rajshahi & 01712000003\\\hline
C04 & Mitu & Khulna & 01712000004\\
\hline

\end{tabular}
\chSixNextColumn
\subsubsection{Account Table}



\small
\setlength{\tabcolsep}{2.5pt}

\begin{tabular}{|l|l|l|l|}
\hline

\chSixHeader
AccountNo & CustomerID & AccountType & Balance\\
\hline
A1001 & C01 & Savings & 150000\\\hline
A1002 & C02 & Current & 85000\\\hline
A1003 & C01 & Current & 70000\\\hline
A1004 & C03 & Savings & 220000\\\hline
A1005 & C04 & Savings & 45000\\
\hline

\end{tabular}
\end{chSixSideBySide}
\begin{chSixSideBySide}[.48]
\subsubsection{Loan Table}



\small

\begin{tabular}{|l|l|l|l|}
\hline

\chSixHeader
LoanID & CustomerID & LoanType & Amount\\
\hline
L01 & C01 & Home & 500000\\\hline
L02 & C02 & Personal & 150000\\\hline
L03 & C03 & Business & 700000\\\hline
L04 & C04 & Personal & 120000\\
\hline

\end{tabular}
\chSixNextColumn
একজন গ্রাহকের একাধিক ব্যাংক হিসাব থাকতে পারে।

তাই ধারণাগতভাবে—

\[
Customer
\quad
1:M
\quad
Account
\]
\end{chSixSideBySide}
\noindent\begin{minipage}{\linewidth}
\chSixExample{৭.৪৯}{সব ব্যাংক হিসাবের তথ্য দেখাও।}

\begin{chSixSideBySide}[.32]
\begin{verbatim}
SELECT *
FROM Account;
\end{verbatim}
\chSixNextColumn
\textbf{ফল}\par\smallskip
{\raggedright\footnotesize\setlength{\tabcolsep}{2pt}
\begin{tabular}{|l|l|l|l|}\hline\chSixHeader
AccountNo & CustomerID & AccountType & Balance\\\hline
A1001 & C01 & Savings & 150000\\\hline
A1002 & C02 & Current & 85000\\\hline
A1003 & C01 & Current & 70000\\\hline
A1004 & C03 & Savings & 220000\\\hline
A1005 & C04 & Savings & 45000\\
\hline\end{tabular}\par}
\end{chSixSideBySide}
\end{minipage}\par\smallskip

\noindent\begin{minipage}{\linewidth}
\chSixExample{৭.৫০}{যেসব হিসাবের Balance 50000-এর বেশি, সেগুলো দেখাও।}

\noindent\begin{minipage}[c]{.46\linewidth}
\begin{verbatim}
SELECT AccountNo, CustomerID, Balance
FROM Account
WHERE Balance>50000;
\end{verbatim}
\end{minipage}\hfill\begin{minipage}[c]{.04\linewidth}\raggedright{\color{chSixTeal}$\Rightarrow$}\end{minipage}\hfill
\begin{minipage}[c]{.47\linewidth}\raggedright\footnotesize\setlength{\tabcolsep}{2pt}\textbf{ফল}\par\smallskip
\begin{tabular}{|l|l|l|}\hline\chSixHeader
AccountNo & CustomerID & Balance\\\hline
A1001 & C01 & 150000\\\hline
A1002 & C02 & 85000\\\hline
A1003 & C01 & 70000\\\hline
A1004 & C03 & 220000\\
\hline\end{tabular}
\end{minipage}
\end{minipage}\par\smallskip

\noindent\begin{minipage}{\linewidth}
\chSixExample{৭.৫১}{ঋণের পরিমাণের সর্বোচ্চ মান বের করো।}

\noindent\begin{minipage}[c]{.53\linewidth}
\begin{verbatim}
SELECT MAX(Amount) AS HighestLoan
FROM Loan;
\end{verbatim}
\end{minipage}\hfill\begin{minipage}[c]{.04\linewidth}\raggedright{\color{chSixTeal}$\Rightarrow$}\end{minipage}\hfill
\begin{minipage}[c]{.40\linewidth}\raggedright\footnotesize\setlength{\tabcolsep}{2pt}\textbf{ফল}\par\smallskip
\begin{tabular}{|l|}\hline\chSixHeader
HighestLoan\\\hline
700000\\
\hline\end{tabular}
\end{minipage}
\end{minipage}\par\smallskip

\noindent\begin{minipage}{\linewidth}
\chSixExample{৭.৫২}{Savings হিসাবগুলোর মধ্যে যেগুলোর Balance 100000-এর বেশি, সেগুলো দেখাও।}

\noindent\begin{minipage}[c]{.46\linewidth}
\begin{verbatim}
SELECT AccountNo, CustomerID, Balance
FROM Account
WHERE AccountType='Savings'
  AND Balance>100000;
\end{verbatim}
\end{minipage}\hfill\begin{minipage}[c]{.04\linewidth}\raggedright{\color{chSixTeal}$\Rightarrow$}\end{minipage}\hfill
\begin{minipage}[c]{.47\linewidth}\raggedright\footnotesize\setlength{\tabcolsep}{2pt}\textbf{ফল}\par\smallskip
\begin{tabular}{|l|l|l|}\hline\chSixHeader
AccountNo & CustomerID & Balance\\\hline
A1001 & C01 & 150000\\\hline
A1004 & C03 & 220000\\
\hline\end{tabular}
\end{minipage}
\end{minipage}\par\smallskip

ফলে A1001 ও A1004 আসবে।

\noindent\begin{minipage}{\linewidth}
\chSixExample{৭.৫৩}{যেসব হিসাবের Balance 80000 থেকে 180000-এর মধ্যে, সেগুলো দেখাও।}

\noindent\begin{minipage}[c]{.53\linewidth}
\begin{verbatim}
SELECT AccountNo, Balance
FROM Account
WHERE Balance BETWEEN 80000 AND 180000;
\end{verbatim}
\end{minipage}\hfill\begin{minipage}[c]{.04\linewidth}\raggedright{\color{chSixTeal}$\Rightarrow$}\end{minipage}\hfill
\begin{minipage}[c]{.40\linewidth}\raggedright\footnotesize\setlength{\tabcolsep}{2pt}\textbf{ফল}\par\smallskip
\begin{tabular}{|l|l|}\hline\chSixHeader
AccountNo & Balance\\\hline
A1001 & 150000\\\hline
A1002 & 85000\\
\hline\end{tabular}
\end{minipage}
\end{minipage}\par\smallskip

\noindent\begin{minipage}{\linewidth}
\chSixExample{৭.৫৪}{সব হিসাবের Balance যোগ করে মোট অর্থের পরিমাণ বের করো।}

\noindent\begin{minipage}[c]{.53\linewidth}
\begin{verbatim}
SELECT SUM(Balance) AS TotalBalance
FROM Account;
\end{verbatim}
\end{minipage}\hfill\begin{minipage}[c]{.04\linewidth}\raggedright{\color{chSixTeal}$\Rightarrow$}\end{minipage}\hfill
\begin{minipage}[c]{.40\linewidth}\raggedright\footnotesize\setlength{\tabcolsep}{2pt}\textbf{ফল}\par\smallskip
\begin{tabular}{|l|}\hline\chSixHeader
TotalBalance\\\hline
570000\\
\hline\end{tabular}
\end{minipage}
\end{minipage}\par\smallskip

\par TotalBalance = 150000+85000+70000+220000+45000 = 570000.\par

\noindent\begin{minipage}{\linewidth}
\chSixExample{৭.৫৫}{Balance-এর গড় বের করো।}

\noindent\begin{minipage}[c]{.53\linewidth}
\begin{verbatim}
SELECT AVG(Balance) AS AverageBalance
FROM Account;
\end{verbatim}
\end{minipage}\hfill\begin{minipage}[c]{.04\linewidth}\raggedright{\color{chSixTeal}$\Rightarrow$}\end{minipage}\hfill
\begin{minipage}[c]{.40\linewidth}\raggedright\footnotesize\setlength{\tabcolsep}{2pt}\textbf{ফল}\par\smallskip
\begin{tabular}{|l|}\hline\chSixHeader
AverageBalance\\\hline
114000\\
\hline\end{tabular}
\end{minipage}
\end{minipage}\par\smallskip

\par AverageBalance = 570000/5 = 114000.\par

\noindent\begin{minipage}{\linewidth}
\chSixExample{৭.৫৬}{Balance বড় থেকে ছোট ক্রমে ব্যাংক হিসাবগুলো সাজাও।}

\noindent\begin{minipage}[c]{.53\linewidth}
\begin{verbatim}
SELECT AccountNo, Balance
FROM Account
ORDER BY Balance DESC;
\end{verbatim}
\end{minipage}\hfill\begin{minipage}[c]{.04\linewidth}\raggedright{\color{chSixTeal}$\Rightarrow$}\end{minipage}\hfill
\begin{minipage}[c]{.40\linewidth}\raggedright\footnotesize\setlength{\tabcolsep}{2pt}\textbf{ফল}\par\smallskip
\chSixResultList{AccountNo, Balance}{\mbox{A1004, 220000}; \mbox{A1001, 150000}; \mbox{A1002, 85000}; \mbox{A1003, 70000}; \mbox{A1005, 45000}}
\end{minipage}
\end{minipage}\par\smallskip

\noindent\begin{minipage}{\linewidth}
\chSixExample{৭.৫৭}{A1002 হিসাবে 15000 টাকা জমা হলে তার Balance হালনাগাদ করো।}

\noindent\begin{minipage}[c]{.53\linewidth}
\begin{verbatim}
UPDATE Account
SET Balance=Balance+15000
WHERE AccountNo='A1002';
\end{verbatim}
\end{minipage}\hfill\begin{minipage}[c]{.04\linewidth}\raggedright{\color{chSixTeal}$\Rightarrow$}\end{minipage}\hfill
\begin{minipage}[c]{.40\linewidth}\raggedright\footnotesize\setlength{\tabcolsep}{2pt}\textbf{ফল}\par\smallskip
\begin{tabular}{|l|l|}\hline\chSixHeader
AccountNo & Balance\\\hline
A1002 & 100000\\
\hline\end{tabular}
\end{minipage}
\end{minipage}\par\smallskip

এই হিসাবে আগে 85000 টাকা ছিল। নতুন করে 15000 টাকা জমা হওয়ায় Balance হলো 85000+15000=100000 টাকা। তাই UPDATE-এ পুরোনো Balance-এর সঙ্গে জমার পরিমাণ যোগ করা হয়েছে।

\begin{keypointbox}{Calculated UPDATE}

\par\noindent\begin{minipage}{\linewidth}
\begin{verbatim}
SET Balance=Balance+15000
\end{verbatim}
\end{minipage}\par

এখানে পুরোনো Balance-এর সঙ্গে নতুন জমার পরিমাণ যোগ করা হচ্ছে।

\end{keypointbox}

\subsection{Salary/Payroll টেবিলে বেতন হিসাব: রাশি ও সারসংক্ষেপের Query}

নিচের সারণিতে HR (Human Resources) মানবসম্পদ বিভাগ এবং IT (Information Technology) তথ্যপ্রযুক্তি বিভাগ বোঝায়।

\begin{flushleft}\begin{minipage}{\linewidth}\raggedright

\scriptsize

\begin{tabular}{|l|l|l|l|l|l|l|}
\hline

\chSixHeader
EmpID & Name & Department & BasicSalary & Allowance & Medical & Deduction\\
\hline
E101 & Rafi & Accounts & 30000 & 5000 & 1500 & 2000\\\hline
E102 & Nila & HR & 35000 & 6000 & 1800 & 2500\\\hline
E103 & Sami & IT & 45000 & 8000 & 2200 & 3000\\\hline
E104 & Mitu & Accounts & 28000 & 4500 & 1400 & 1500\\\hline
E105 & Hasan & IT & 50000 & 9000 & 2500 & 3500\\
\hline

\end{tabular}

\captionof{table}{Salary Table}

\end{minipage}\end{flushleft}

\chSixExample{৭.৫৮}{প্রত্যেক কর্মীর প্রাপ্য বেতন বা Net Salary দেখাও।}

ধরি—

\par NetSalary = BasicSalary + Allowance + Medical - Deduction.\par

SQL—

\noindent\begin{minipage}{\linewidth}
\noindent\begin{minipage}[c]{.53\linewidth}
\begin{verbatim}
SELECT Name, BasicSalary + Allowance
       + Medical - Deduction AS NetSalary
FROM Salary;
\end{verbatim}
\end{minipage}\hfill\begin{minipage}[c]{.04\linewidth}\raggedright{\color{chSixTeal}$\Rightarrow$}\end{minipage}\hfill
\begin{minipage}[c]{.40\linewidth}\raggedright\footnotesize\setlength{\tabcolsep}{2pt}\textbf{ফল}\par\smallskip
\chSixResultList{Name, NetSalary}{\mbox{Rafi, 34500}; \mbox{Nila, 40300}; \mbox{Sami, 52200}; \mbox{Mitu, 32400}; \mbox{Hasan, 58000}}
\end{minipage}
\end{minipage}\par\smallskip

এখানে NetSalary কুয়েরির মাধ্যমে হিসাব করা ফল। মূল টেবিলে এই নামে আলাদা ফিল্ড থাকা বাধ্যতামূলক নয়।

\noindent\begin{minipage}{\linewidth}
\chSixExample{৭.৫৯}{যেসব কর্মীর Basic Salary 35000-এর বেশি, তাঁদের তথ্য দেখাও।}

\noindent\begin{minipage}[c]{.53\linewidth}
\begin{verbatim}
SELECT EmpID, Name, BasicSalary
FROM Salary
WHERE BasicSalary>35000;
\end{verbatim}
\end{minipage}\hfill\begin{minipage}[c]{.04\linewidth}\raggedright{\color{chSixTeal}$\Rightarrow$}\end{minipage}\hfill
\begin{minipage}[c]{.40\linewidth}\raggedright\footnotesize\setlength{\tabcolsep}{2pt}\textbf{ফল}\par\smallskip
\begin{tabular}{|l|l|l|}\hline\chSixHeader
EmpID & Name & BasicSalary\\\hline
E103 & Sami & 45000\\\hline
E105 & Hasan & 50000\\
\hline\end{tabular}
\end{minipage}
\end{minipage}\par\smallskip

ফলাফলে Sami ও Hasan-এর তথ্য পাওয়া যাবে।

\noindent\begin{minipage}{\linewidth}
\chSixExample{৭.৬০}{IT বিভাগের কর্মীদের তথ্য দেখাও।}

\noindent\begin{minipage}[c]{.53\linewidth}
\begin{verbatim}
SELECT Name, BasicSalary
FROM Salary
WHERE Department='IT';
\end{verbatim}
\end{minipage}\hfill\begin{minipage}[c]{.04\linewidth}\raggedright{\color{chSixTeal}$\Rightarrow$}\end{minipage}\hfill
\begin{minipage}[c]{.40\linewidth}\raggedright\footnotesize\setlength{\tabcolsep}{2pt}\textbf{ফল}\par\smallskip
\begin{tabular}{|l|l|}\hline\chSixHeader
Name & BasicSalary\\\hline
Sami & 45000\\\hline
Hasan & 50000\\
\hline\end{tabular}
\end{minipage}
\end{minipage}\par\smallskip

\noindent\begin{minipage}{\linewidth}
\chSixExample{৭.৬১}{Basic Salary-এর সর্বোচ্চ মান বের করো।}

\noindent\begin{minipage}[c]{.53\linewidth}
\begin{verbatim}
SELECT MAX(BasicSalary) AS HighestSalary
FROM Salary;
\end{verbatim}
\end{minipage}\hfill\begin{minipage}[c]{.04\linewidth}\raggedright{\color{chSixTeal}$\Rightarrow$}\end{minipage}\hfill
\begin{minipage}[c]{.40\linewidth}\raggedright\footnotesize\setlength{\tabcolsep}{2pt}\textbf{ফল}\par\smallskip
\begin{tabular}{|l|}\hline\chSixHeader
HighestSalary\\\hline
50000\\
\hline\end{tabular}
\end{minipage}
\end{minipage}\par\smallskip

\noindent\begin{minipage}{\linewidth}
\chSixExample{৭.৬২}{Basic Salary-এর গড় বের করো।}

\noindent\begin{minipage}[c]{.53\linewidth}
\begin{verbatim}
SELECT AVG(BasicSalary) AS AverageSalary
FROM Salary;
\end{verbatim}
\end{minipage}\hfill\begin{minipage}[c]{.04\linewidth}\raggedright{\color{chSixTeal}$\Rightarrow$}\end{minipage}\hfill
\begin{minipage}[c]{.40\linewidth}\raggedright\footnotesize\setlength{\tabcolsep}{2pt}\textbf{ফল}\par\smallskip
\begin{tabular}{|l|}\hline\chSixHeader
AverageSalary\\\hline
37600\\
\hline\end{tabular}
\end{minipage}
\end{minipage}\par\smallskip

\noindent\begin{minipage}{\linewidth}
\chSixExample{৭.৬৩}{প্রতিটি বিভাগের গড় Basic Salary দেখাও।}

\noindent\begin{minipage}[c]{.53\linewidth}
\begin{verbatim}
SELECT Department,
       AVG(BasicSalary) AS AverageSalary
FROM Salary
GROUP BY Department;
\end{verbatim}
\end{minipage}\hfill\begin{minipage}[c]{.04\linewidth}\raggedright{\color{chSixTeal}$\Rightarrow$}\end{minipage}\hfill
\begin{minipage}[c]{.40\linewidth}\raggedright\footnotesize\setlength{\tabcolsep}{2pt}\textbf{ফল}\par\smallskip
\begin{tabular}{|l|l|}\hline\chSixHeader
Department & AverageSalary\\\hline
Accounts & 29000\\\hline
HR & 35000\\\hline
IT & 47500\\
\hline\end{tabular}
\end{minipage}
\end{minipage}\par\smallskip

\noindent\begin{minipage}{\linewidth}
\chSixExample{৭.৬৪}{যেসব বিভাগের গড় Basic Salary 30000-এর বেশি, সেগুলো দেখাও।}

\noindent\begin{minipage}[c]{.53\linewidth}
\begin{verbatim}
SELECT Department,
       AVG(BasicSalary) AS AverageSalary
FROM Salary
GROUP BY Department
HAVING AVG(BasicSalary)>30000;
\end{verbatim}
\end{minipage}\hfill\begin{minipage}[c]{.04\linewidth}\raggedright{\color{chSixTeal}$\Rightarrow$}\end{minipage}\hfill
\begin{minipage}[c]{.40\linewidth}\raggedright\footnotesize\setlength{\tabcolsep}{2pt}\textbf{ফল}\par\smallskip
\begin{tabular}{|l|l|}\hline\chSixHeader
Department & AverageSalary\\\hline
HR & 35000\\\hline
IT & 47500\\
\hline\end{tabular}
\end{minipage}
\end{minipage}\par\smallskip

HR ও IT বিভাগের গড় বেতন শর্ত পূরণ করে, তাই এই বিভাগ দুটির তথ্য পাওয়া যাবে।

\Needspace{18\baselineskip}
\subsection{লাইব্রেরির ডেটাবেজে Query}

\subsubsection{Book Table}

\begin{flushleft}\begin{minipage}{\linewidth}\raggedright

\small

\begin{tabular}{|l|l|l|l|l|}
\hline

\chSixHeader
BookID & BookName & Author & Price & Category\\
\hline
B01 & Database Basics & Rahman & 450 & Computer\\\hline
B02 & Easy SQL & Karim & 520 & Computer\\\hline
B03 & Physics Plus & Nabil & 600 & Science\\\hline
B04 & Web Design & Rafi & 480 & Computer\\\hline
B05 & English Grammar & Hasan & 350 & Language\\
\hline

\end{tabular}

\captionof{table}{Book Table}

\end{minipage}\end{flushleft}

\noindent\begin{minipage}{\linewidth}
\chSixExample{৭.৬৫}{Computer শ্রেণির বইগুলোর তথ্য দেখাও।}

\noindent\begin{minipage}[c]{.46\linewidth}
\begin{verbatim}
SELECT BookID, BookName, Author
FROM Book
WHERE Category='Computer';
\end{verbatim}
\end{minipage}\hfill\begin{minipage}[c]{.04\linewidth}\raggedright{\color{chSixTeal}$\Rightarrow$}\end{minipage}\hfill
\begin{minipage}[c]{.47\linewidth}\raggedright\footnotesize\setlength{\tabcolsep}{2pt}\textbf{ফল}\par\smallskip
\begin{tabular}{|l|l|l|}\hline\chSixHeader
BookID & BookName & Author\\\hline
B01 & Database Basics & Rahman\\\hline
B02 & Easy SQL & Karim\\\hline
B04 & Web Design & Rafi\\
\hline\end{tabular}
\end{minipage}
\end{minipage}\par\smallskip

\noindent\begin{minipage}{\linewidth}
\chSixExample{৭.৬৬}{যেসব বইয়ের Price 500-এর বেশি, সেগুলোর তথ্য দেখাও।}

\noindent\begin{minipage}[c]{.53\linewidth}
\begin{verbatim}
SELECT BookName, Price
FROM Book
WHERE Price>500;
\end{verbatim}
\end{minipage}\hfill\begin{minipage}[c]{.04\linewidth}\raggedright{\color{chSixTeal}$\Rightarrow$}\end{minipage}\hfill
\begin{minipage}[c]{.40\linewidth}\raggedright\footnotesize\setlength{\tabcolsep}{2pt}\textbf{ফল}\par\smallskip
\begin{tabular}{|l|l|}\hline\chSixHeader
BookName & Price\\\hline
Easy SQL & 520\\\hline
Physics Plus & 600\\
\hline\end{tabular}
\end{minipage}
\end{minipage}\par\smallskip

\noindent\begin{minipage}{\linewidth}
\chSixExample{৭.৬৭}{Rahman-এর লেখা বইগুলোর তথ্য দেখাও।}

\noindent\begin{minipage}[c]{.53\linewidth}
\begin{verbatim}
SELECT BookName, Price
FROM Book
WHERE Author='Rahman';
\end{verbatim}
\end{minipage}\hfill\begin{minipage}[c]{.04\linewidth}\raggedright{\color{chSixTeal}$\Rightarrow$}\end{minipage}\hfill
\begin{minipage}[c]{.40\linewidth}\raggedright\footnotesize\setlength{\tabcolsep}{2pt}\textbf{ফল}\par\smallskip
\begin{tabular}{|l|l|}\hline\chSixHeader
BookName & Price\\\hline
Database Basics & 450\\
\hline\end{tabular}
\end{minipage}
\end{minipage}\par\smallskip

\subsubsection{LibraryIssue Table}

\begin{flushleft}\begin{minipage}{\linewidth}\raggedright

\small

\begin{tabular}{|l|l|l|l|l|}
\hline

\chSixHeader
IssueID & StudentID & BookID & IssueDate & ReturnDate\\
\hline
I01 & S01 & B01 & 2026-08-01 & 2026-08-07\\\hline
I02 & S02 & B03 & 2026-08-03 & NULL\\\hline
I03 & S01 & B04 & 2026-08-04 & NULL\\\hline
I04 & S05 & B02 & 2026-08-05 & 2026-08-12\\
\hline

\end{tabular}

\captionof{table}{LibraryIssue Table}

\end{minipage}\end{flushleft}

\noindent\begin{minipage}{\linewidth}
\chSixExample{৭.৬৮}{যেসব বই এখনো ফেরত দেওয়া হয়নি, সেগুলোর তথ্য দেখাও।}

\noindent\begin{minipage}[c]{.53\linewidth}
\begin{verbatim}
SELECT StudentID, BookID
FROM LibraryIssue
WHERE ReturnDate IS NULL;
\end{verbatim}
\end{minipage}\hfill\begin{minipage}[c]{.04\linewidth}\raggedright{\color{chSixTeal}$\Rightarrow$}\end{minipage}\hfill
\begin{minipage}[c]{.40\linewidth}\raggedright\footnotesize\setlength{\tabcolsep}{2pt}\textbf{ফল}\par\smallskip
\begin{tabular}{|l|l|}\hline\chSixHeader
StudentID & BookID\\\hline
S02 & B03\\\hline
S01 & B04\\
\hline\end{tabular}
\end{minipage}
\end{minipage}\par\smallskip

B03 ও B04 বইয়ের ফেরতের তারিখ নেই, তাই এগুলো ফলাফলে দেখা যাবে।

\noindent\begin{minipage}{\linewidth}
\chSixExample{৭.৬৯}{ধার দেওয়া বইগুলোর নামসহ সংশ্লিষ্ট তথ্য দেখাও।}

\noindent\begin{minipage}[c]{.46\linewidth}
\begin{verbatim}
SELECT i.IssueID,
       b.BookName,
       i.StudentID
FROM LibraryIssue AS i
INNER JOIN Book AS b
ON i.BookID=b.BookID;
\end{verbatim}
\end{minipage}\hfill\begin{minipage}[c]{.04\linewidth}\raggedright{\color{chSixTeal}$\Rightarrow$}\end{minipage}\hfill
\begin{minipage}[c]{.47\linewidth}\raggedright\footnotesize\setlength{\tabcolsep}{2pt}\textbf{ফল}\par\smallskip
\begin{tabular}{|l|l|l|}\hline\chSixHeader
IssueID & BookName & StudentID\\\hline
I01 & Database Basics & S01\\\hline
I02 & Physics Plus & S02\\\hline
I03 & Web Design & S01\\\hline
I04 & Easy SQL & S05\\
\hline\end{tabular}
\end{minipage}
\end{minipage}\par\smallskip

এখানে— \texttt{LibraryIssue.BookID} এবং \texttt{Book.BookID} মিলিয়ে সম্পর্কিত রেকর্ড একসঙ্গে দেখানো হয়েছে।

\subsection{Sales টেবিলে গাণিতিক রাশি ও সমষ্টিগত হিসাব}

\begin{flushleft}\begin{minipage}{\linewidth}\raggedright

\small

\begin{tabular}{|l|l|l|l|l|}
\hline

\chSixHeader
SaleID & Product & UnitPrice & Quantity & SalesPerson\\
\hline
S01 & Keyboard & 1200 & 2 & Rafi\\\hline
S02 & Mouse & 650 & 5 & Nila\\\hline
S03 & Headphone & 1800 & 3 & Rafi\\\hline
S04 & PenDrive & 900 & 4 & Sami\\\hline
S05 & Speaker & 2200 & 2 & Nila\\
\hline

\end{tabular}

\captionof{table}{Sales Table}

\end{minipage}\end{flushleft}

\noindent\begin{minipage}{\linewidth}
\chSixExample{৭.৭০}{প্রতিটি বিক্রয়ের TotalPrice দেখাও।}

\noindent\begin{minipage}[c]{.53\linewidth}
\begin{verbatim}
SELECT Product,
       UnitPrice*Quantity AS TotalPrice
FROM Sales;
\end{verbatim}
\end{minipage}\hfill\begin{minipage}[c]{.04\linewidth}\raggedright{\color{chSixTeal}$\Rightarrow$}\end{minipage}\hfill
\begin{minipage}[c]{.40\linewidth}\raggedright\footnotesize\setlength{\tabcolsep}{2pt}\textbf{ফল}\par\smallskip
\chSixResultList{Product, TotalPrice}{\mbox{Keyboard, 2400}; \mbox{Mouse, 3250}; \mbox{Headphone, 5400}; \mbox{PenDrive, 3600}; \mbox{Speaker, 4400}}
\end{minipage}
\end{minipage}\par\smallskip

\noindent\begin{minipage}{\linewidth}
\chSixExample{৭.৭১}{সব বিক্রয়ের মূল্য যোগ করে মোট বিক্রয়মূল্য বের করো।}

\noindent\begin{minipage}[c]{.53\linewidth}
\begin{verbatim}
SELECT SUM(UnitPrice*Quantity)
       AS TotalSales
FROM Sales;
\end{verbatim}
\end{minipage}\hfill\begin{minipage}[c]{.04\linewidth}\raggedright{\color{chSixTeal}$\Rightarrow$}\end{minipage}\hfill
\begin{minipage}[c]{.40\linewidth}\raggedright\footnotesize\setlength{\tabcolsep}{2pt}\textbf{ফল}\par\smallskip
\begin{tabular}{|l|}\hline\chSixHeader
TotalSales\\\hline
19050\\
\hline\end{tabular}
\end{minipage}
\end{minipage}\par\smallskip

\noindent\begin{minipage}{\linewidth}
\chSixExample{৭.৭২}{প্রত্যেক SalesPerson-এর মোট বিক্রয়মূল্য দেখাও।}

\noindent\begin{minipage}[c]{.53\linewidth}
\begin{verbatim}
SELECT SalesPerson,
       SUM(UnitPrice*Quantity) AS TotalSales
FROM Sales
GROUP BY SalesPerson;
\end{verbatim}
\end{minipage}\hfill\begin{minipage}[c]{.04\linewidth}\raggedright{\color{chSixTeal}$\Rightarrow$}\end{minipage}\hfill
\begin{minipage}[c]{.40\linewidth}\raggedright\footnotesize\setlength{\tabcolsep}{2pt}\textbf{ফল}\par\smallskip
\begin{tabular}{|l|l|}\hline\chSixHeader
SalesPerson & TotalSales\\\hline
Nila & 7650\\\hline
Rafi & 7800\\\hline
Sami & 3600\\
\hline\end{tabular}
\end{minipage}
\end{minipage}\par\smallskip

